% =====================================================================================================
\section{Threat Model and Security Goals}\label{sec:threat}
\textit{System.} A camera node (the PYNQ-Z2 board) encodes video and sends one JPEG per frame over UDP to a viewer; the viewer returns heartbeats and alerts over a feedback channel. Both ends hold a 128-bit master key $K_{\mathrm{m}}$ provisioned at installation. As in SRTP, whose key management is a separate protocol~\cite{rfc3711}, and in every scheme we compare with~\cite{vandroogenbroeck2002,kishore2015,auer2013,hirose2024,bagbaba2015}, key distribution is out of scope; a standard key exchange on the ARM core at session start is the natural extension (Section~\ref{sec:discussion}).

\textit{Adversary.} The attacker sits on the network path. It observes all traffic and may modify, drop, delay, replay, reorder or inject both video and feedback datagrams; it does not hold the keys. It knows the design, the code and the format, may know some frames in clear (known plaintext), may choose content that the camera films (chosen plaintext), and may train neural networks on any amount of public video. A passive eavesdropper sends nothing, so it cannot be detected; against it only the configured minimum protection level helps. Physical and side-channel attacks on the device are out of scope; a deployment that needs them can use a threshold implementation of Ascon, which Gro{\ss} et al.\ built at about 3.1 times the area of their unprotected low-area core~\cite{gross2017}.

\textit{Assets and goals.}
\begin{itemize}
\item[G1] \emph{Content confidentiality:} a keyless observer should not recover the shape of the scene. Format compliance means that the frame size, the image dimensions, the protection level and coarse statistics are public by design; we measure what else leaks (Section~\ref{sec:security}) and what the public parts reveal (Section~\ref{sec:discussion}).
\item[G2] \emph{Integrity and authenticity} of every frame together with its metadata (level, session, counter, size): any change must be rejected.
\item[G3] \emph{Freshness:} no replayed or re-ordered old frame is shown as current.
\item[G4] \emph{Availability signal:} the viewer must learn that the camera is alive, and the camera that the viewer receives; silence must be treated as an attack.
\item[G5] \emph{Forward secrecy across epochs:} a video key captured later does not decrypt earlier epochs.
\end{itemize}
G2--G5 are enforced cryptographically (Section~\ref{sec:design}); G1 is a property of the protection level, which we measure empirically against the strongest attackers we could build. The receiver accepts a frame only if its tag verifies, so a wrong key, a modified frame or a forged frame is rejected rather than shown as noise; a viewer without the key can still play the stream, which is the point of format compliance.
